El resultado más importante está en la pendiente de la ecuación \eqref{eq:ajuste}. La teoría pide que valga exactamente uno, y salió 1.005, apenas medio por ciento más.

Eso quiere decir algo muy concreto.

Por cada centímetro que alejábamos el objeto del espejo, la imagen también se alejaba un centímetro del otro lado, sin importar si estábamos a 8 o a 25 cm, y la relación entre las dos distancias se mantuvo igual de principio a fin, que es justo lo que dice la ecuación \eqref{eq:distancias}.

La ordenada al origen, en cambio, no fue cero sino $-0.32$ cm. Y aquí aparece un detalle interesante: en las once posiciones la imagen quedó del lado de menor distancia, nunca del otro, como también se ve en la tabla \ref{tab:distancias}, donde todas las diferencias son positivas. Cuando un error se repite siempre con el mismo signo no se debe a la mala suerte, sino a algo del montaje. Lo más probable es que la línea de referencia no coincidiera exactamente con la superficie reflejante del espejo, ya sea porque el espejo no quedó justo encima de ella o porque en muchos espejos la capa que refleja está detrás del vidrio, unos milímetros atrás de la cara frontal. Un corrimiento de unos tres milímetros en esa posición basta para explicar la diferencia promedio de 0.24 cm.

También influyen otras cosas, aunque menos. El grosor del lápiz y de las líneas, la dificultad de juzgar a ojo el momento exacto en que desaparece el paralaje y el hecho de que la regla solo marca milímetros pueden mover cada medida uno o dos milímetros, y como ese error pesa más cuando la distancia es chica, la posición 3, a solo 8 cm, fue la que tuvo el error relativo más alto, de 5 por ciento.

En conclusión, comprobamos que en un espejo plano la imagen se forma detrás del espejo, sobre la perpendicular que pasa por el objeto y a la misma distancia que este, con una pendiente experimental de 1.005, un ajuste de $R^2 = 0.9996$ y un error relativo promedio de 1.8 por ciento. Las pequeñas diferencias que aparecieron se explican sobre todo por un corrimiento constante en la posición del espejo.

El método de paralaje resultó sencillo y sorprendentemente preciso.
